% 4_resultados_esperados.tex
\chapter{RESULTADOS ESPERADOS}

Ao final deste projeto, espera-se alcançar uma compreensão aprofundada das capacidades fotocatalíticas dos tungstatos metálicos estudados na geração de ROS e de seu potencial em aplicações bactericidas. Serão desenvolvidos modelos computacionais baseados em termodinâmica ab initio para descrever as principais terminações superficiais desses materiais, possibilitando a identificação das configurações mais estáveis. As propriedades eletrônicas serão investigadas por meio de cálculos de estrutura de bandas, DOS e função trabalho, permitindo avaliar o potencial fotocatalítico dos sistemas analisados. Além disso, será estudada a interação das moléculas de O$_2$ e H$_2$O com as superfícies mais estáveis, com o objetivo de elucidar os mecanismos de ativação envolvidos na formação de ROS, incluindo a identificação de sítios nucleofílicos e eletrofílicos relevantes para os processos de adsorção. Os cálculos de energia de adsorção fornecerão informações quantitativas sobre a afinidade das diferentes superfícies pelas moléculas alvo. Complementarmente, a investigação dos estados de transição, obtidos por meio de cálculos utilizando o NEB e validados por análises de frequências de fônons, permitirá determinar as barreiras energéticas associadas às reações, oferecendo uma visão abrangente da reatividade superficial dos materiais. Por fim, o projeto contribuirá de forma significativa para a formação do doutorando, fortalecendo sua trajetória acadêmica, bem como a produção científica e a colaboração entre os grupos e instituições envolvidos.


